\section{Functional CALM: Proofs}
\label{app:proofs-sec2}

This appendix proves the results of Section~\ref{sec:functional-calm}.

\begin{proof}[Proof of Lemma~\ref{lem:run-residual}]
The run starts at the canonical state $\sigma_P(i_0)$, whose table is
$P[\mathrm{in}\InputGe i_0]$. Each input event $\mathsf{input}(j)$ keeps only
the rows with $\mathrm{in}\InputGe j$ (Equation~\ref{eq:input-event}), and each
output event $\mathsf{output}(o)$ keeps only the rows with
$\mathrm{out}\OutputGe o$ (Equation~\ref{eq:output-action}). The input
conditions reduce to the last one, $\mathrm{in}\InputGe i$, because each later
input is an extension of the earlier ones. One output condition remains for each
emitted output.
\end{proof}

\begin{proof}[Proof of Proposition~\ref{prop:confluence}]
A common state with current input $j$ has a row at $j$, and since rows are
never added, that row belongs to every $R_r$.  Conversely, let $(j,u)$ belong
to every $R_r$.  Then $\mathsf{input}(j)$ is defined at every $(i_r,R_r)$, and
$\mathsf{output}(u)$ is defined at $(j,R_r[\mathrm{in}\InputGe j])$, since
$u\in R_r(j)$.  By
Lemma~\ref{lem:run-residual}, $(j,u)\in R_r$ says that $u$ \outrel every output
emitted on run $r$.  After $\mathsf{input}(j)$ and $\mathsf{output}(u)$, the
same lemma leaves exactly
$P[\mathrm{in}\InputGe j,\ \mathrm{out}\OutputGe u]$, independent of the
starting run.
\end{proof}

\begin{proof}[Proof of Lemma~\ref{lem:confluence-legality}]
We use legality in a form that covers both functional and nondeterministic
specs: $o$ is legal at $(i,R)$ when, for every row $(k,x)\in R$, some
$u\in R(k)$ \outrel both $o$ and $x$.  In a functional table the only row at
$k$ is $(k,x)$, so this says $x$ \outrel $o$
(Definition~\ref{def:exposable}); with several outputs per input it says that
$\Excluded_o$ is empty (Section~\ref{sec:lifted-output-events}).  An output
event $\mathsf{output}(o)$ from $(i,R)$ applies a \commitment exactly when $o$
is not legal at $(i,R)$, so the lemma's condition says that every output a run
can emit is legal at the state where it is emitted.

($\Rightarrow$) Suppose an output event of some run applies a \commitment: the
run ends at $(i,R)$ before it, and $o\in R(i)$ is not legal there.  Then some $(k,x)\in R$ has no $u\in R(k)$ \outrelpart both $o$ and
$x$.  We extend the run in two ways: by $\mathsf{output}(o)$, and by
$\mathsf{input}(k)$ followed by $\mathsf{output}(x)$, which is defined since
$(k,x)$ remains after $\mathsf{input}(k)$.  The input $k$ is a common \indesc
of the two current inputs.  A row at $k$ in both resulting tables would be a
row $(k,u)\in R$ with $u$ \outrelpart both $o$ and $x$, and there is none.  By
Proposition~\ref{prop:confluence}, the two runs are not confluent at $k$.

($\Leftarrow$) Suppose no output event applies a \commitment, so every output
that a run can emit is legal, and let runs
end at $(i_r,R_r)$, $r=1,\ldots,m$, with common \indesc $j$.  Each $R_r$ keeps a
row at its current input, say $(i_r,u_r)$, and by Lemma~\ref{lem:run-residual}
$u_r$ \outrel every output emitted on run $r$.  The canonical state
$\sigma_P(i_r)$ is reached by the empty run and has $u_r$ as an output, so $u_r$
is legal there.  By totality there is some $v_0\in P(j)$.  For $r=1,\ldots,m$
in turn, $(j,v_{r-1})$ is a row of $\sigma_P(i_r)$, since $j\InputGe i_r$, so
legality of $u_r$ gives $v_r\in P(j)$ \outrelpart both $u_r$ and $v_{r-1}$.
Then $v_m$ \outrel every $u_r$.  By Lemma~\ref{lem:run-residual},
$(j,v_m)\in R_r$ for every $r$, since $j\InputGe i_r$ and $v_m$ \outrel every
output emitted on run $r$.  Proposition~\ref{prop:confluence} gives
confluence at $j$.
\end{proof}

\begin{proof}[Proof of Corollary~\ref{cor:state-by-input}]
By Lemma~\ref{lem:run-residual}, the run's remainder table is
$P[\mathrm{in}\InputGe i,\ \mathrm{out}\OutputGe o_1,\ \ldots,\
\mathrm{out}\OutputGe o_m]$, where each $o_\ell=f(i_\ell)$ for some
$i_\ell\InputOrd i$.  Every row whose input \inrel $i$ has the form $(k,f(k))$ with
$i_\ell\InputOrd k$, so monotonicity gives $o_\ell\OutputOrd f(k)$, and the
output conditions remove nothing.
\end{proof}

\begin{proof}[Proof of Proposition~\ref{prop:functional-exposure-boundary}]
Suppose $f$ is monotone, and let a \commitment-free run reach $(i,R)$.  By Corollary~\ref{cor:state-by-input},
$(i,R)=\sigma_P(i)$, so the legal $j\InputGe i$ are exactly the \indescs of
$i$, and monotonicity gives $f(i)\OutputOrd f(j)$ at each, so $f(i)$ is a
legal output (Definition~\ref{def:exposable}).
Conversely, if every such output is a legal output, then at $\sigma_P(i)$,
which the empty run reaches, legality of $f(i)$ gives $f(i)\OutputOrd f(j)$
for every \indesc $j$ of $i$, so $f$ is monotone.
\end{proof}

\begin{proof}[Proof of Theorem~\ref{thm:functional-calm}]
By Lemma~\ref{lem:confluence-legality}, the runs are consistent exactly when
no output event of a run applies a \commitment, that is, when every output that
a run can emit is legal.  This condition is equivalent to
legality at every state reachable without \commitment.  One direction is
immediate, since such states are reached by runs.  For the other, suppose every
output at every state reachable without \commitment is legal.  Then, by
induction along any run, each of its output events is legal, so the run is
\commitment-free and every state it reaches is reachable without \commitment;
hence every output it can emit is legal.  By
Proposition~\ref{prop:functional-exposure-boundary}, legality at every state
reachable without \commitment holds exactly when $f$ is monotone.
\end{proof}
